\chapter{Conclusiones}

A lo largo de este trabajo se ha abordado el proceso de análisis, diseño, desarrollo y puesta en marcha de una aplicación de software. El primer paso fue analizar los problemas existentes en las plataformas de video actuales con el objetivo de crear una alternativa que optimice las oportunidades de mejora identificadas. A partir del diagnóstico inicial, se diseñó y desarrolló una solución tecnológica orientada a resolver las necesidades de los usuarios, explorando distintos patrones arquitectónicos y herramientas disponible en \href{http://boostervideos.net}{\underline{www.boostervideos.net}}. 

Durante la ejecución del proyecto, se analizaron diferentes proveedores, herramientas, servicios y tecnologías disponibles para implementar una aplicación de streaming de video. Esto permitió comparar alternativas para procesos clave como el almacenamiento de video, la transcodificación y segmentación, la distribución de contenido, la autenticación, el despliegue y la lógica de negocio. Asimismo, se evaluaron las ventajas y desventajas de cada opción teniendo en cuenta los costos, sus limitaciones y su grado de complejidad. También se presentaron casos de uso públicos en los que estas tecnologías son utilizadas por empresas o plataformas disponibles en el mercado. Esto brindó una visión más clara de la utilidad de las herramientas demostrando que el proyecto incorpora criterios de análisis propios de un entorno profesional de desarrollo de software. Esta etapa fue fundamental para comprender que el desarrollo de una aplicación de esta naturaleza, además de escribir código, implica analizar servicios de terceros, evaluar alternativas y tomar decisiones de diseño coherentes con los objetivos del proyecto.

La puesta en marcha del servicio web de la plataforma también fue un aspecto importante del trabajo. Esto significó afrontar desafíos en un entorno de producción, resolver problemas de configuración, gestionar errores, integrar servicios de terceros, controlar la estabilidad del sistema y monitorear los recursos utilizados por la aplicación. Además, se realizaron actividades publicitarias creando un video  \href{http://boostervideos.net}{\underline{anuncio}}  y realizando publicaciones en redes sociales para atraer usuarios. Esto, a su vez, implicó una observación cuidadosa de los recursos consumidos por el servicio. Con un incremento en el volumen de usuarios, era necesario vigilar los límites para garantizar la continuidad del sistema y evitar caídas abruptas. Fue particularmente relevante, el intercambio con usuarios reales permitiendo corregir y adaptar funcionalidades derivadas de sugerencias específicas.

Otro aspecto fundamental del proyecto fue la identificación de vulnerabilidades. Durante el desarrollo se detectaron problemas relacionados con la seguridad que afectaban seriamente la calidad del servicio. Además, se introdujeron límites de uso para aliviar la carga sobre el sistema y evitar un uso indebido de los recursos. Estas vulnerabilidades fueron analizadas y mitigadas mediante mejoras en el código y cambios en la arquitectura.

Por último, el proyecto ha permitido experimentar de forma completa el proceso de creación de una aplicación de software y de un proyecto empresarial. Se tomaron decisiones de negocio orientadas a la optimización de costos y se realizaron campañas publicitarias en redes sociales para captar 300 usuarios activos. Además, se enfrentaron desafíos tecnológicos y problemas técnicos para desarrollar y mantener una plataforma en Internet.

Como resultado, el trabajo demuestra el proceso de construcción de una solución tecnológica a partir de la identificación y resolución de problemas, aplicando conceptos técnicos y estratégicos para crear una plataforma de streaming de video alternativa y mejorada respecto de las opciones disponibles en el mercado.


